\section{Verificación}

El repositorio tiene \textbf{\totalpruebas{} pruebas automáticas}, organizadas como un
espejo del código, que cubren el \coberturapruebas{}\,\% de las líneas y ramas del motor,
y pasa limpio el análisis de estilo (\texttt{ruff}) y el de tipos en modo estricto
(\texttt{mypy}). Tres ideas las guían:

\begin{itemize}[nosep, leftmargin=*]
  \item \textbf{Validadores estructurales.} Una estructura puede responder bien y estar
        rota por dentro. Cada índice tiene un validador que la recorre entera: en el B+, el
        orden, la ocupación mínima y que todas las hojas estén a la misma profundidad; en
        el hash, que cada cubeta reciba $2^{gd-ld}$ punteros; en el R-Tree, que cada MBR
        sea exactamente el ajustado.
  \item \textbf{Comparación con un modelo.} Cada estructura se somete a miles de
        operaciones al azar y se compara con una lista ordenada o con la fuerza bruta. Cada
        consulta espacial se ejecuta sin índice y con R-Tree, y cada condición sobre una
        tabla sin índices, un heap con índices, un secuencial y un B+ agrupado: las
        respuestas tienen que coincidir.
  \item \textbf{Contraste externo.} Las 1\,800 consultas del experimento espacial devuelven
        exactamente las mismas filas que PostGIS.
\end{itemize}

Las pruebas de almacenamiento usan páginas de 256 bytes a propósito: así cualquier prueba
cruza varias páginas y los errores de frontera aparecen enseguida.
